\paragrafos{Επίκεντρη γωνία, τόξο και μέτρηση τόξου}{Β.1.12, «Επίκεντρη γωνία – Σχέση επίκεντρης γωνίας και του αντίστοιχου τόξου – Μέτρηση τόξου»}

\begin{anagke}
Στο πάρκο δίπλα στο σχολείο υπάρχει μια ρόδα λούνα παρκ με $8$ ίσες καμπίνες. Η παρέα την κοιτάζει από ψηλά: οι καμπίνες είναι στερεωμένες σε ίσες αποστάσεις πάνω σε ένα κύκλο. «Πόσες μοίρες γυρίζει ο τροχός για να πάει η μία καμπίνα στη θέση της επόμενης;» ρωτάει ο Ηλίας. Η Πολυξένη σκέφτεται: «Ο πλήρης κύκλος είναι $360^\circ$ και υπάρχουν $8$ ίσα διαστήματα: $360:8=45^\circ$».

Η Αγγελική σχεδιάζει δύο ακτίνες προς δύο διαδοχικές καμπίνες και παρατηρεί τη γωνία ανάμεσά τους με κορυφή το κέντρο: «Είναι $45^\circ$, και το κομμάτι του κύκλου ανάμεσα στις δύο καμπίνες το μετράμε και αυτό $45^\circ$». Ο Στέλιος ρωτάει: «Το μήκος του κύκλου και οι μοίρες είναι το ίδιο;». Η κυρία Αλεξάνδρα απαντάει: «Όχι. Άλλο το μήκος σε μέτρα και άλλο το άνοιγμα σε μοίρες. Σήμερα μιλάμε για το άνοιγμα, το λεγόμενο \emph{μέτρο του τόξου}».

Στην παράγραφο αυτή ορίζουμε την επίκεντρη γωνία και το αντίστοιχο τόξο, και μετράμε τα τόξα με μοίρες.
\end{anagke}

\begin{lexi}{επίκεντρη γωνία, τόξο}{επίκεντρη γωνία, τόξο}
Η \lexiword{επίκεντρη} γωνία είναι η γωνία που έχει την κορυφή της «επί του κέντρου» του κύκλου. Το \lexiword{τόξο} είναι το κομμάτι του κύκλου, σαν το τόξο του τοξοβόλου, καμπυλωμένο. Το \lexiword{ημικύκλιο} είναι το μισό του κύκλου.
\end{lexi}

\begin{metafrasi}
Γωνία με κορυφή το κέντρο. & επίκεντρη γωνία $\widehat{\mathrm{ΑΟΒ}}$\\
Το κομμάτι του κύκλου μέσα στη γωνία. & αντίστοιχο τόξο $\overset{\frown}{\mathrm{ΑΒ}}$\\
Το μέτρο του τόξου. & ίσο με το μέτρο της επίκεντρης γωνίας\\
Ο πλήρης κύκλος. & $360^\circ$\\
Το ημικύκλιο. & $180^\circ$\\
\end{metafrasi}

\tmima{Η θεωρία}

\begin{orismos}{επίκεντρη γωνία και τόξο}
\textbf{Επίκεντρη} γωνία ενός κύκλου $(\mathrm{Ο},\rho)$ είναι κάθε γωνία $x\mathrm{Ο}y$ με κορυφή το κέντρο $\mathrm{Ο}$. Αν οι πλευρές της τέμνουν τον κύκλο στα $\mathrm{Α}$ και $\mathrm{Β}$, το τμήμα του κύκλου που βρίσκεται στο εσωτερικό της γωνίας λέγεται \textbf{αντίστοιχο τόξο} της και γράφεται $\overset{\frown}{\mathrm{ΑΒ}}$.
\end{orismos}

\begin{center}
\begin{tikzpicture}[scale=1.2]
\draw (0,0) circle (1.6);
\fill (0,0) circle (2pt) node[below left] {$\mathrm{Ο}$};
\draw[thick] (0,0)--(1.6,0) (0,0)--({1.6*cos(50)},{1.6*sin(50)});
\fill (1.6,0) circle (2pt) node[right] {$\mathrm{Α}$}; \fill ({1.6*cos(50)},{1.6*sin(50)}) circle (2pt) node[above right] {$\mathrm{Β}$};
\draw[gblue,very thick] (1.6,0) arc(0:50:1.6);
\draw[gblue] (0.5,0) arc(0:50:0.5); \node[gblue] at (0.8,0.2) {\small$\omega$};
\end{tikzpicture}
\end{center}

\begin{kanonas}{μέτρο τόξου}
Το \textbf{μέτρο ενός τόξου} σε μοίρες είναι ίσο με το μέτρο της αντίστοιχης επίκεντρης γωνίας. Ολόκληρος ο κύκλος είναι τόξο $360^\circ$ και το ημικύκλιο $180^\circ$. Ίσες επίκεντρες γωνίες αντιστοιχούν σε ίσα τόξα, και αντίστροφα.
\end{kanonas}

Έτσι το τόξο $\overset{\frown}{\mathrm{ΑΒ}}$ του σχήματος είναι $50^\circ$. Όταν η γωνία είναι μη κυρτή, το τόξο που αντιστοιχεί σε αυτήν είναι το μεγάλο: $360^\circ-50^\circ=310^\circ$.

Το μέτρο τόξου (σε μοίρες) είναι διαφορετικό από το \emph{μήκος} του τόξου (σε μέτρα): δύο ομόκεντροι κύκλοι έχουν τόξα με ίδιο μέτρο αλλά διαφορετικό μήκος.

\begin{prosoxi}
\begin{itemize}[leftmargin=1.2em,itemsep=0.3ex,topsep=2pt]
\item Το μέτρο του τόξου σε μοίρες δεν είναι το μήκος του σε μέτρα.
\item Κάθε χορδή $\mathrm{ΑΒ}$ ορίζει δύο τόξα (ένα μικρό και ένα μεγάλο), εκτός αν είναι διάμετρος (δύο ημικύκλια).
\item Η επίκεντρη γωνία έχει την κορυφή της στο κέντρο: αν η κορυφή είναι αλλού, δεν είναι επίκεντρη.
\item Σε διαφορετικούς κύκλους, ίσα τόξα (σε μοίρες) δεν σημαίνει ίσα μήκη.
\end{itemize}
\end{prosoxi}

\tmima{Λυμένα παραδείγματα}

\paradeigma{Η ρόδα με τις $8$ καμπίνες}
\tixero Οκτώ ίσες καμπίνες πάνω σε κύκλο.\par
\tithelo Το μέτρο του τόξου ανάμεσα σε δύο διαδοχικές.\par
\possindeo Ο κύκλος είναι $360^\circ$ και τα $8$ τόξα είναι ίσα: $360^\circ:8=45^\circ$.\par
\apantisi $45^\circ$.\par
\elegxos $8\cdot45=360$.

\paradeigma{Τόξο και επίκεντρη}
\tixero Η επίκεντρη γωνία $\widehat{\mathrm{ΑΟΒ}}=50^\circ$.\par
\tithelo Το μικρό και το μεγάλο τόξο $\overset{\frown}{\mathrm{ΑΒ}}$.\par
\possindeo Το μικρό τόξο είναι $50^\circ$. Το μεγάλο είναι $360^\circ-50^\circ=310^\circ$.\par
\apantisi $50^\circ$ και $310^\circ$.\par
\elegxos $50+310=360$.

\paradeigma{Διάμετρος και τόξα}
\tixero Η $\mathrm{ΑΒ}$ είναι διάμετρος του κύκλου.\par
\tithelo Το μέτρο των τόξων που ορίζει.\par
\possindeo Η διάμετρος χωρίζει τον κύκλο σε δύο ημικύκλια. Κάθε ημικύκλιο είναι $180^\circ$.\par
\apantisi $180^\circ$ το καθένα.\par
\elegxos Η επίκεντρη γωνία $\widehat{\mathrm{ΑΟΒ}}$ είναι ευθεία.

\paradeigma{Ένα ρολόι}
\tixero Ο λεπτοδείκτης κινείται από το $12$ στο $4$.\par
\tithelo Το τόξο που διέγραψε η άκρη του.\par
\possindeo Κάθε λεπτό ($1/60$ του κύκλου) είναι $6^\circ$, ή κάθε πεντάλεπτο $30^\circ$. Από το $12$ στο $4$ είναι $4$ πεντάλεπτα, $4\cdot30^\circ=120^\circ$.\par
\apantisi $120^\circ$.\par
\elegxos Στα $20$ λεπτά, $20\cdot6^\circ=120^\circ$.

\tmima{Ασκήσεις}

\zoni{Το κατάλαβα}{τόξα και μοίρες}

\ex Σχεδίασε κύκλο και επίκεντρη γωνία $70^\circ$. Ποιο είναι το μέτρο του μικρού και του μεγάλου τόξου που αντιστοιχούν;

\ex Ένας κύκλος χωρίζεται σε $5$ ίσα τόξα. Πόσες μοίρες είναι το καθένα; Και σε $9$; Σε $12$;

\ex Πόσες μοίρες είναι το $\frac{1}{4}$, το $\frac{1}{8}$, το $\frac{3}{4}$ και το $\frac{5}{6}$ του κύκλου;

\ex Ένα τόξο έχει μέτρο $60^\circ$. Τι μέρος του κύκλου είναι;

\ex Μια διάμετρος $\mathrm{ΑΒ}$ ορίζει δύο τόξα. Πόσες μοίρες είναι το καθένα;

\ex Σχεδίασε κύκλο και τρεις επίκεντρες διαδοχικές γωνίες $100^\circ$, $120^\circ$, $140^\circ$. Πόσο είναι τα τόξα και πόσο το άθροισμά τους;

\ex Σωστό ή λάθος; α) Το μέτρο του τόξου είναι ίσο με της επίκεντρης γωνίας. \quad β) Το ημικύκλιο είναι $180^\circ$. \quad γ) Ίσα τόξα έχουν πάντα ίσα μήκη σε οποιουσδήποτε κύκλους. \quad δ) Ο κύκλος είναι τόξο $360^\circ$.

\zoni{Το χρησιμοποιώ}{ρόδες, ρολόγια, πίτες}

\ex \textbf{Η ρόδα.} Μια ρόδα λούνα παρκ έχει $12$ ίσες καμπίνες. Πόσες μοίρες γυρίζει για να έρθει η επόμενη καμπίνα στη θέση της;

\ex \textbf{Το ρολόι.} Πόσες μοίρες διαγράφει η άκρη του λεπτοδείκτη σε $20$ λεπτά; Σε $45$ λεπτά; Σε μία ώρα;

\ex \textbf{Η πίτα.} Μια πίτα κόβεται σε φέτες με επίκεντρες γωνίες $45^\circ$. Πόσες φέτες βγαίνουν; Τι μέρος της πίτας είναι $3$ φέτες;

\ex \textbf{Το ραντάρ.} Ένα ραντάρ περιστρέφεται και σαρώνει ένα τόξο $75^\circ$. Τι μέρος του κύκλου καλύπτει;

\ex \textbf{Ο ανεμόμυλος.} Ένας ανεμόμυλος με $4$ πτερύγια γυρίζει $90^\circ$. Πόσα τέτοια τόξα χρειάζονται για μία πλήρη περιστροφή;

\zoni{Το σκέφτομαι}{μέτρο ή μήκος;}

\ex Δύο κύκλοι έχουν ίδιο κέντρο και διαφορετικές ακτίνες. Οι ίδιες δύο ακτίνες ορίζουν στον καθένα ένα τόξο. Είναι ίσα τα μέτρα τους; Είναι ίσα τα μήκη τους;

\ex Γιατί ένας πλήρης κύκλος είναι $360^\circ$ και όχι $100^\circ$; Θα άλλαζε τίποτα αν το διαλέγαμε $100$;

\ex Μια χορδή $\mathrm{ΑΒ}$ ορίζει τόξο $120^\circ$. Πόσο είναι το άλλο τόξο της;

\ex Αν η επίκεντρη γωνία διπλασιαστεί, τι γίνεται το τόξο;

\zoni{Για το σπίτι}{ο κύκλος και οι μοίρες}

\ex Σχεδίασε έναν κύκλο και χώρισέ τον σε $6$ ίσα τόξα με διαβήτη και μοιρογνωμόνιο. Πόσες μοίρες είναι το καθένα;

\ex Βρες κάτι στο σπίτι (ρολόι, πίτα, ρόδα) που χωρίζεται σε ίσα τόξα και υπολόγισε τις μοίρες κάθε τόξου.
